\section{Bối cảnh và tính cấp thiết của đề tài}

Trong bối cảnh chuyển đổi số diễn ra mạnh mẽ, dữ liệu ngày càng trở thành nguồn lực quan trọng đối với hoạt động quản lý và ra quyết định trong doanh nghiệp. Đặc biệt trong lĩnh vực kinh doanh bán lẻ, dữ liệu phát sinh từ doanh số bán hàng, thông tin khách hàng, sản phẩm và khu vực phân phối ngày càng lớn, đòi hỏi phải có các phương pháp phân tích phù hợp để khai thác hiệu quả. Việc chỉ quan sát số liệu theo cách thủ công hoặc dựa trên kinh nghiệm cảm tính không còn đáp ứng tốt yêu cầu quản trị trong môi trường cạnh tranh hiện nay.

Bên cạnh đó, doanh nghiệp không chỉ quan tâm đến mức doanh thu đạt được mà còn cần hiểu rõ xu hướng tiêu dùng của khách hàng. Thông qua việc phân tích dữ liệu bán hàng, doanh nghiệp có thể nhận diện sự biến động doanh số theo thời gian, đánh giá hiệu quả của từng khu vực, nhóm sản phẩm và phân khúc khách hàng, đồng thời xác định các yếu tố có ảnh hưởng đến lợi nhuận và hiệu quả kinh doanh. Đây là cơ sở quan trọng để hỗ trợ nhà quản lý đưa ra các quyết định phù hợp hơn trong hoạt động bán hàng và phân bổ nguồn lực.

Xuất phát từ yêu cầu thực tiễn đó, đề tài ``Phân tích doanh số và xu hướng tiêu dùng khách hàng'' được em lựa chọn nhằm nghiên cứu dữ liệu bán hàng dưới góc độ phân tích dữ liệu hiện đại. Đề tài tập trung vào việc làm sạch dữ liệu, phân tích khám phá dữ liệu, trực quan hóa kết quả và áp dụng một số mô hình học máy để đánh giá tình hình kinh doanh cũng như nhận diện xu hướng tiêu dùng của khách hàng.

Tính cấp thiết của đề tài thể hiện ở chỗ việc khai thác dữ liệu bán hàng một cách có hệ thống không chỉ giúp nâng cao tính khách quan trong đánh giá hoạt động kinh doanh mà còn góp phần hỗ trợ doanh nghiệp đưa ra quyết định dựa trên dữ liệu. Đồng thời, đề tài cũng có ý nghĩa thực tiễn đối với quá trình học tập và nghiên cứu, vì cho phép vận dụng các kiến thức về phân tích dữ liệu, trực quan hóa dữ liệu và học máy vào một bài toán cụ thể trong lĩnh vực kinh doanh.

\section{Mục tiêu nghiên cứu}

\textbf{Mục tiêu tổng quát:}

Phân tích doanh số bán hàng và xu hướng tiêu dùng của khách hàng dựa trên bộ dữ liệu Superstore Sales nhằm đánh giá tình hình hoạt động kinh doanh, làm rõ các yếu tố ảnh hưởng đến doanh thu và lợi nhuận, đồng thời hỗ trợ rút ra các nhận xét có cơ sở dữ liệu phục vụ cho việc đánh giá hiệu quả bán hàng và hỗ trợ ra quyết định.

\textbf{Mục tiêu cụ thể:}

\begin{itemize}
    \item Hệ thống hóa cơ sở lý thuyết về phân tích dữ liệu trong kinh doanh, trực quan hóa dữ liệu, học máy và các chỉ số đánh giá mô hình nhằm làm nền tảng cho quá trình triển khai đề tài.
    
    \item Thực hiện quá trình đọc dữ liệu, kiểm tra dữ liệu, làm sạch, chuẩn hóa và xử lý các vấn đề như dữ liệu thiếu, dữ liệu trùng lặp, dữ liệu ngoại lai trên môi trường \textbf{Python}, qua đó bảo đảm bộ dữ liệu đạt yêu cầu cho các bước phân tích tiếp theo.
    
    \item Tiến hành phân tích tổng quan và phân tích chuyên sâu bộ dữ liệu theo các khía cạnh như thời gian, khu vực, danh mục sản phẩm, tiểu danh mục sản phẩm, phân khúc khách hàng, doanh thu, chiết khấu và lợi nhuận nhằm làm rõ đặc điểm hoạt động kinh doanh và xu hướng tiêu dùng của khách hàng.
    
    \item Thiết kế dashboard trực quan trên \textbf{Power BI} nhằm thể hiện các chỉ số và biểu đồ quan trọng, từ đó hỗ trợ quá trình theo dõi, so sánh và đánh giá tình hình kinh doanh dưới nhiều góc độ khác nhau.
    
    \item Xây dựng và đánh giá các mô hình gồm \textbf{Hồi quy tuyến tính đa biến (Multiple Linear Regression)}, \textbf{Random Forest Regressor}, \textbf{Gradient Boosting Regressor}, \textbf{XGBoost Regressor} và \textbf{K-Means Clustering} nhằm phục vụ cho bài toán dự báo, phân cụm khách hàng, so sánh hiệu quả giữa các phương pháp và rút ra các nhận xét tổng hợp về doanh số bán hàng cũng như xu hướng tiêu dùng của khách hàng.
\end{itemize}


\section{Đối tượng và phạm vi nghiên cứu}

\textbf{Đối tượng nghiên cứu:}

Đối tượng nghiên cứu của đề tài là dữ liệu bán hàng trong bộ dữ liệu \textbf{Superstore Sales}, bao gồm các thông tin liên quan đến đơn hàng, khách hàng, sản phẩm, khu vực bán hàng, doanh thu, số lượng bán, chiết khấu và lợi nhuận. Trên cơ sở đó, đề tài tập trung nghiên cứu mối quan hệ giữa các yếu tố trong dữ liệu nhằm làm rõ đặc điểm hoạt động kinh doanh và xu hướng tiêu dùng của khách hàng.

\textbf{Phạm vi nghiên cứu:}

\begin{itemize}
    \item Về nội dung, đề tài tập trung vào việc phân tích doanh số bán hàng, doanh thu, lợi nhuận và xu hướng tiêu dùng của khách hàng thông qua các khía cạnh như thời gian, khu vực, danh mục sản phẩm, tiểu danh mục sản phẩm và phân khúc khách hàng.
    
    \item Về dữ liệu, đề tài sử dụng bộ dữ liệu \textbf{Superstore Sales} làm nguồn dữ liệu nghiên cứu chính, đồng thời thực hiện các bước tiền xử lý, trực quan hóa và xây dựng mô hình trên bộ dữ liệu này.
    
    \item Về phương pháp, đề tài triển khai phân tích dữ liệu trên môi trường \textbf{Python}, xây dựng dashboard trực quan trên \textbf{Power BI} và áp dụng một số mô hình như \textbf{Hồi quy tuyến tính đa biến}, \textbf{Random Forest Regressor}, \textbf{Gradient Boosting Regressor}, \textbf{XGBoost Regressor} và \textbf{K-Means Clustering} để phục vụ cho quá trình phân tích, dự báo và so sánh.
    
    \item Về mục đích, đề tài giới hạn trong việc đánh giá tình hình kinh doanh, nhận diện xu hướng tiêu dùng và xem xét các yếu tố ảnh hưởng đến doanh thu, lợi nhuận; không đi sâu vào việc xây dựng hệ thống triển khai thực tế trong doanh nghiệp.
\end{itemize}


\section{Phương pháp nghiên cứu}

Để thực hiện đề tài \textit{``Phân tích doanh số và xu hướng tiêu dùng khách hàng''}, quá trình nghiên cứu được triển khai theo hướng kết hợp giữa phân tích dữ liệu, trực quan hóa dữ liệu và áp dụng các mô hình học máy. Trước hết, đề tài tiến hành thu thập và khai thác bộ dữ liệu \textbf{Superstore Sales} làm nguồn dữ liệu chính cho toàn bộ quá trình nghiên cứu. Trên cơ sở đó, dữ liệu được đưa vào môi trường \textbf{Python} để kiểm tra, làm sạch, chuẩn hóa và xử lý các vấn đề như dữ liệu thiếu, dữ liệu trùng lặp và giá trị ngoại lai nhằm bảo đảm chất lượng dữ liệu đầu vào.

Sau bước tiền xử lý, đề tài thực hiện phân tích dữ liệu theo hai mức độ. Thứ nhất là phân tích tổng quan nhằm nhận diện đặc điểm chung của bộ dữ liệu, bao gồm cấu trúc dữ liệu, các biến nghiên cứu và phân bố ban đầu của một số chỉ tiêu quan trọng. Thứ hai là phân tích chuyên sâu theo các khía cạnh như thời gian, khu vực, danh mục sản phẩm, tiểu danh mục sản phẩm và phân khúc khách hàng nhằm làm rõ tình hình doanh số, doanh thu, lợi nhuận cũng như xu hướng tiêu dùng của khách hàng.

Bên cạnh đó, đề tài kết hợp \textbf{Power BI} để xây dựng dashboard trực quan, thể hiện kết quả phân tích thông qua các biểu đồ, thẻ chỉ số và bộ lọc tương tác. Việc trực quan hóa dữ liệu giúp quá trình quan sát, so sánh và đánh giá kết quả trở nên rõ ràng hơn, đồng thời hỗ trợ tốt cho việc trình bày và rút ra nhận xét.

Trên cơ sở dữ liệu đã được xử lý và phân tích, đề tài tiếp tục triển khai một số mô hình trên môi trường \textbf{Python} gồm \textbf{Hồi quy tuyến tính đa biến (Multiple Linear Regression)}, \textbf{Random Forest Regressor}, \textbf{Gradient Boosting Regressor}, \textbf{XGBoost Regressor} và \textbf{K-Means Clustering}. Các mô hình này được sử dụng nhằm phục vụ cho việc dự báo, phân tích mối quan hệ giữa các biến và phân cụm khách hàng theo đặc điểm tiêu dùng. Kết quả của các mô hình được đánh giá thông qua các chỉ số phù hợp như \textbf{MAE}, \textbf{MSE}, \textbf{RMSE}, \textbf{\(R^2\)}, đồng thời được so sánh để xác định phương pháp phù hợp hơn đối với bộ dữ liệu nghiên cứu.

Cuối cùng, từ kết quả phân tích dữ liệu, dashboard trực quan và các mô hình đã xây dựng, đề tài tiến hành tổng hợp, đánh giá và rút ra các nhận xét về tình hình kinh doanh, xu hướng tiêu dùng của khách hàng và các yếu tố ảnh hưởng đến hiệu quả bán hàng. Đây là cơ sở để hình thành kết luận của đề tài và làm rõ giá trị ứng dụng của phân tích dữ liệu trong hỗ trợ ra quyết định kinh doanh.



\section{Ý nghĩa khoa học và thực tiễn của đề tài}

\textbf{Ý nghĩa khoa học:}

Đề tài góp phần hệ thống hóa và vận dụng các kiến thức về phân tích dữ liệu trong kinh doanh, trực quan hóa dữ liệu và học máy vào một bài toán cụ thể trong lĩnh vực bán hàng. Thông qua việc nghiên cứu bộ dữ liệu Superstore Sales, đề tài làm rõ quy trình xử lý dữ liệu, phân tích dữ liệu và đánh giá mô hình trong bối cảnh dữ liệu kinh doanh. Bên cạnh đó, việc kết hợp giữa phân tích khám phá dữ liệu, xây dựng dashboard trực quan và triển khai nhiều mô hình khác nhau cũng giúp làm rõ khả năng ứng dụng của các phương pháp phân tích dữ liệu hiện đại trong việc khai thác thông tin từ dữ liệu bán hàng.

Ngoài ra, đề tài còn có ý nghĩa về mặt phương pháp khi thực hiện so sánh giữa các mô hình như Hồi quy tuyến tính đa biến, Random Forest Regressor, Gradient Boosting Regressor, XGBoost Regressor và K-Means Clustering trên cùng một bộ dữ liệu nghiên cứu. Điều này giúp làm rõ đặc điểm, hiệu quả và mức độ phù hợp của từng phương pháp trong bài toán phân tích doanh số và xu hướng tiêu dùng khách hàng.

\textbf{Ý nghĩa thực tiễn:}

Về mặt thực tiễn, đề tài giúp đánh giá tình hình hoạt động kinh doanh thông qua các chỉ tiêu như doanh số, doanh thu, lợi nhuận, khu vực bán hàng, danh mục sản phẩm và phân khúc khách hàng. Kết quả phân tích cho phép nhận diện xu hướng tiêu dùng của khách hàng, xác định các nhóm sản phẩm hoặc khu vực hoạt động hiệu quả, đồng thời làm rõ ảnh hưởng của một số yếu tố như chiết khấu, số lượng bán hay phân khúc khách hàng đến kết quả kinh doanh.

Bên cạnh đó, dashboard trực quan được xây dựng trong đề tài có ý nghĩa hỗ trợ quan sát, theo dõi và so sánh dữ liệu một cách rõ ràng, trực quan hơn. Các mô hình phân tích dữ liệu được triển khai cũng góp phần nâng cao tính khách quan trong việc đánh giá và dự báo, từ đó hỗ trợ quá trình ra quyết định dựa trên dữ liệu. Mặc dù đề tài được thực hiện trên bộ dữ liệu nghiên cứu, quy trình phân tích và phương pháp triển khai vẫn có giá trị tham khảo cho các bài toán phân tích kinh doanh tương tự trong thực tế.